\section*{Required author Section 5.3: recursive factorization}
This formalizes and extends recursive factorization, giving notions extendable to matrix products (Definitions5.13/5.15), a simplification (Lemma5.25) and refinement (Proposition5.24) for later use.
\paragraph{Definition 5.9.} Let $L$ be the $\binom{[n]}{d/2}\times\binom{[n]}{\le d/2}$ matrix
\[L(I,A):=\sum_{\substack{T':|V(T')\cup I\cup A|\le\tau,\ A=S_l(I,A;T')\\T'\cap E(A)=\emptyset,\ (I,A;T')\text{ left-generated}}}(\omega/n)^{|V(T')\cup I\cup A|-|A|}\chi_{T'}.\tag{5.19}\]
Let $Q_0$ be the $\binom{[n]}{\le d/2}\times\binom{[n]}{\le d/2}$ matrix
\[Q_0(A,B):=\sum_{\substack{T_m:|T_m\cup A\cup B|\le\tau\\A,B\in\mathrm{mSep}_{A,B}(T_m)}}(\omega/n)^{|V(T_m)\cup A\cup B|}\chi_{T_m}.\tag{5.20}\]
Finally, $D:=\operatorname{diag}((\omega/n)^{|A|/2})_{A\in\binom{[n]}{\le d/2}}$ (5.21). The author calls $L(DQ_0)L^\top$ the first-approximate diagonalization of $M'$. Despite the name, $M'-L(DQ_0D)L^\top$ (5.22) is far from negligible. Recursive factorization will in the end give
\[M'=L[D(Q_0-Q_1+Q_2\cdots\pm Q_{d/2})D]L^\top+E\tag{5.23}\]
for negligible error-matrix $E$.
\paragraph{Remark 5.10.} Use of $D$ is superficial; it keeps middle matrices $Q_i$ better-positioned. $LD$ corresponds to the ``$L$'' matrix of [5].
\paragraph{Definition 5.11 ([5],Definition6.5).} The canonical decomposition of ribbon $R=(I,J;T)$ is the uniquely determined ribbon-triple
\[(R_l,R_m,R_r)=((I,A;T_l),(A,B;T_m),(B,J;T_r)),\]
where $A=S_l(I,J;T)$, $B=S_r(I,J;T)$. $V(R_l)$ is $A$ unioned with vertices reachable from $I$ in $T$ without touching $A$, and $T_l=T|_{V(R_l)}\setminus E(A)$. $V(R_r)$ is symmetric, reachable from $J$ without touching $B$, with $T_r=T|_{V(R_r)}\setminus E(B)$. Finally $T_m=T\setminus(T'\sqcup T'')$. These are the left, middle and right ribbon of $R$.
\paragraph{Remark 5.12.} The canonical decomposition satisfies
\begin{gather*}A=S_l(I,A;T_l),\quad B=S_r(B,J;T_r),\\T_l\cap E(A)=\emptyset=T_r\cap E[A],\\R_l\text{ left-generated},\quad R_r\text{ right-generated},\\A,B\in\mathrm{mSep}_{A,B}(T_m)\quad(\text{so }|A|=|B|).
\end{gather*}
These are the inner conditions. The outer intersection property is
\[V(R_l)\cap V(R_m)\subseteq A,\quad V(R_m)\cap V(R_r)\subseteq B,\quad V(R_l)\cap V(R_r)\subseteq A\cap B,\]
implying $e(R_l)+|V(R_m)|+e(R_r)=|V(R)|$ (5.24).
\paragraph{Definition 5.13 (Inner and outer canonicality).} For a triple of ribbons $(R_l,R_m,R_r)=((I,A;T_l),(A,B;T_m),(B,J;T_r))$ ($T_l,T_m,T_r$ arbitrary subsets of an edge-set), their ribbon-sum is $(I,J;T)$ where $T=T_l\oplus T_m\oplus T_r$. It is inner-canonical if it satisfies the four inner conditions above (5.25). It is outer-canonical if
\[V(R_l)\cap V(R_m)\subseteq A,\quad V(R_m)\cap V(R_r)\subseteq B,\quad V(R_l)\cap V(R_r)\subseteq A\cap B.\tag{5.26}\]
It is canonical if both inner- and outer-canonical.
\paragraph{Proposition 5.14.} Canonical triples are 1-1 correspondent to their ribbon-sum, via canonical decomposition.
\begin{proof}This follows by an immediate check from the definition.\end{proof}
Recall $\mathbb R[\{\chi_T\}]$ is the ring adding fresh variables $\chi_T$ into $\mathbb R$ for each $T\subseteq\binom{[n]}2$, with relations $\chi_{T'}\chi_{T''}=\chi_T$ whenever $T'\oplus T''=T$.
\paragraph{Definition 5.15 (Approximate form).} Suppose matrices $X,Y$ have rows/columns indexed by subsets of $[n]$, entries in $\mathbb R[\{\chi_T\}]$, and every character viewed as a ribbon on (row,column) has size $\le\tau$. Suppose $XYX^\top$ is defined. Every nonzero triple product, without collecting like terms, is
\[X(I,A;T_l)Y(A,B;T_m)X(J,B;T_r)\chi_{T_l}\chi_{T_m}\chi_{T_r}\tag{5.28}\]
in $XYX^\top$ (5.27), with the coefficient nonzero in $\mathbb R$. Identify it with the natural ribbon-triple; its resulting term is $X(I,A;T_l)Y(A,B;T_m)X(J,B;T_r)\chi_{T_l\oplus T_m\oplus T_r}$. Call it outer-canonical if that triple is outer-canonical; it exceeds degree if $|V(T)\cup I\cup J|>\tau$. The approximation form is
\begin{align*}XYX^\top&=(XYX^\top)_{\mathrm{can}}+(XYX^\top)_{\mathrm{non\text{-}can}}+E_{\mathrm{deg}},\\(XYX^\top)_{\mathrm{can}}&=XYX^\top-(XYX^\top)_{\mathrm{non\text{-}can}}-E_{\mathrm{deg}}.\tag{5.29}\end{align*}
Here the canonical part collects outer-canonical products not exceeding degree, the non-canonical part collects non-outer-canonical products, and $E_{\mathrm{deg}}$ all remaining terms.
\paragraph{Remark 5.16.} Proposition5.14 explains Definition5.9 a posteriori: $M'=[L(DQ_0D)L^\top]_{\mathrm{can}}$.
\paragraph{Definition 5.17 (First-approximate factorization of $M'$).}
\[M'=L(DQ_0D)L^\top-[L(DQ_0D)L^\top]_{\mathrm{non\text{-}can}}-E_{1;\mathrm{deg}}.\tag{5.30}\]
$E_{1;\mathrm{deg}}$ is Definition5.15 applied to this product. $L(DQ_0D)L^\top$ is the first-approximate factorization. The high-degree error is negligible in norm (proved analogously in the exact case); matrices considered are supported on clique rows/columns of $G$. The main error is the non-canonical part. The key is that it too approximately factors through $L,L^\top$: there is $Q_1$ with
\[[L(DQ_0D)L^\top]_{\mathrm{non\text{-}can}}=[L(DQ_1D)L^\top]_{\mathrm{can}}+E'_{1;\mathrm{negl}}.\tag{5.31}\]
Repeat for $[L(DQ_1D)L^\top]_{\mathrm{non\text{-}can}}$ and so on.
\paragraph{Definition 5.18 ([5],Definition6.9).} A generalized ribbon is a usual ribbon with a new isolated vertex-set. Write $R^*=(A,B;T^*)$ with $T^*=T\sqcup I$, $T$ an edge-set, $I$ disjoint from $V(T)\cup A\cup B$, denoted $I(R^*)$. $V(R^*)=V(T)\cup A\cup B\cup I$. Usual ribbons have $I=\emptyset$. $(A,B;T)$ is the unique largest ribbon in $R$.
\paragraph{Remark 5.19.} $I(R^*)$ differs from the isolated set of the underlying graph since it excludes $A\cup B$.
\paragraph{Definition 5.20.} A side-inner-canonical triple $((I,A;T_l),(A,B;T_m),(B,J;T_r))$ has left/right ribbons satisfying their inner-canonical conditions (first three of (5.25)); the middle is just a ribbon.
\paragraph{Definition 5.21 (Separating factorization; [5],Definition6.10).} Given a side-inner-canonical triple $(R_l,R_m,R_r)$ as above, put $T=T_l\oplus T_m\oplus T_r$. Let $Z$ be the multiset of unexpected intersections, the vertices from $(R_l\cap R_m)-A$, $(R_m\cap R_r)-B$, $(R_l\cap R_r)-(A\cap B)$. Its intersection size is $z(R_l,R_m,R_r)=|Z|$. Then
\[|V(R_l)\cup V(R_m)\cup V(R_r)|=|V(R_l)|+|V(R_m)|+|V(R_r)|-|A|-|B|-z.\tag{5.32}\]
Define $S'_l$ to be the leftmost min-separator of $(I,A\cup(Z\cap V(R_l));T_l)$, and $S'_r$ the rightmost min-separator of $(B\cup(Z\cap V(R_r)),J;T_r)$. They lie in $V(T)\cup I\cup J$. Define $R'_l=(I,S'_l;T'_l)$: $V(R'_l)$ is $S'_l$ unioned with vertices in $R_l$ reachable from $I$ by paths in $T_l$ without touching $S'_l$, and $T'_l$ is $T_l\setminus E(S'_l)$ restricted to $V(R'_l)$. Define $R'_r$ symmetrically. In particular $T'_l\cap T'_r=\emptyset$. The generalized middle ribbon is $R^*_m=(S'_l,S'_r;T^*_m)$ with
\[T^*_m=T\setminus(T'_l\sqcup T'_r)\sqcup I(R^*_m),\]
where $I(R^*_m)$ collects all remaining isolated vertices:
\[I(R^*_m)=V(R_l)\cup V(R_m)\cup V(R_r)-V(T)\cup I\cup J.\tag{5.33}\]
This separating factorization is denoted $(R_l,R_m,R_r)\to(R'_l,R^*_m,R'_r)$ (5.34).
\paragraph{Remark 5.22.} (1) The resulting triple is side-inner-canonical and outer-canonical. Its pairwise intersections lie within $S'_l,S'_r,S'_l\cap S'_r$. Replacing $R^*_m$ by its largest ribbon gives a canonical triple disjoint from $I(R^*_m)$. (2) $R'_l\subseteq R_l$, and $S'_l$ separates $(V(R'_l),V(R_l)-V(R'_l))$ in $R_l$. The part of $R_l$ to the right of $S'_l$ is disjoint from $R'_l$ and is in $R^*_m$; symmetrically for $R_r$. (3) Since $S'_l$ separates $(I,A)$ in $R_l$ and $A$ is its unique min-separator, there are $|A|$ vertex-disjoint paths from $A$ to $S'_l$ in $R_l$; symmetrically for $R_r$.
